\chapter{Conclusion and Future Work}
\label{Chapter5}

\section{Conclusion}

This thesis set out to design, implement, and evaluate uMatter, a mental-health support platform for teenagers that unifies self-care tracking with professional therapy access, addressing a gap identified between existing self-care applications (Daylio, Bearable) and existing therapy platforms (BetterHelp, Talkspace, Wysa), none of which combine both halves of the care journey behind a single, permission-based data-sharing model.

The resulting system is a seven-service microservices architecture behind an API gateway, with a database-per-service data layer, stateless JWT authentication, and an AI companion grounded in each user's own tracking data. A structured security-hardening process identified and resolved concrete issues across business logic, access control, and data integrity, including a database-level fix for a race condition in appointment booking. The full stack was deployed to cloud infrastructure - migrated from Oracle Cloud to Microsoft Azure over the course of the project - and evaluated both functionally and against the competitive landscape it was designed to improve on.

Taken together, this work shows that a single, security-conscious microservices platform can plausibly serve both sides of adolescent mental health support - self-care and professional therapy - rather than requiring a teenager to use two disconnected products.

% =====================================================================
% NEW SECTION (skeleton) - các feature được đánh giá cao bởi người dùng
% và chuyên gia. Draft tiếng Anh + placeholder; về nhà viết lại kỹ hơn.
% =====================================================================
\section{Features Most Valued by Users and Experts}

Across the user trials and the three expert consultations conducted during this project (Section~\ref{sec:expert-findings}), two parts of uMatter stood out as the most consistently and highly valued: the \textbf{Treasure Box (Hộp Trân Quý)} and the suite of \textbf{daily reflective tracking} features. Both are described in the product tour in Appendix~\ref{...} % TODO: trỏ tới appendix/showcase phù hợp
and are singled out here because they attracted the strongest positive response.

\subsection{The Treasure Box (Hộp Trân Quý)}

% TODO: viết lại thành đoạn văn hoàn chỉnh. Ý chính:
\begin{itemize}
\item \textbf{What it is.} A personal store of positive ``core memories'' - photos, kind messages, small wins, life goals, and motivating reflections - that a user saves ahead of time and reopens whenever they need comforting.
\item \textbf{Why it was valued.} On a low day, negative bias makes good memories hard to recall; a pre-saved, reliable anchor back to the good gives users self-compassion they can reach for exactly when they need it.
\item \textbf{Where it came from.} The feature grew directly out of expert collaboration: Miss Pham Thi Thanh Qui first suggested a feature to reconnect users with their reasons for living and their motivation when feeling low, and mister Vuong Nguyen Toan Thien, the CEO of LUMOS, concretized that idea into the Treasure Box, grounding it in psychological research on positive memory and mood.
\item \textbf{How it fits.} It is one of the destinations in uMatter's ``psychology first-aid'' response to a negative mood, alongside guided breathing and the SOS crisis contacts.
\item \textbf{[placeholder]} % TODO: chèn phản hồi/đánh giá cụ thể của người dùng thử (trích dẫn, số liệu nếu có).
\end{itemize}

\subsection{Daily Reflective Tracking}

% TODO: viết lại thành đoạn văn hoàn chỉnh. Ý chính:
\begin{itemize}
\item \textbf{What it is.} The everyday self-check-in habit at the heart of the app: mood check-ins, the emotional journal, sleep, nutrition and water, automatic step counting, and guided breathing - each surfaced as a friendly ``today'' card with weekly trends.
\item \textbf{Why it was valued.} Docter Pham Tran Thanh Nghiep assessed the daily reflective tracking features as ``already very good'' and already helpful to users' mental wellbeing. Miss Qui emphasised that physical and mental health are inseparable, endorsing the physical-health tracking (steps, sleep, nutrition, breathing) as legitimate mental-health signal.
\item \textbf{What makes it stick.} Gentle Focus-Mode reminders roughly every ninety minutes, daily achievement ``trophies'' for completing a chain of self-care tasks, and streaks turn tracking into a low-friction daily habit rather than a chore.
\item \textbf{Why it matters system-wide.} This daily data is also the foundation that grounds the AI companion, the therapist's view (under permission), and the trusted-friend view - so the feature most valued on its own is also what makes the rest of the platform work.
\item \textbf{[placeholder]} % TODO: chèn phản hồi/đánh giá cụ thể của người dùng thử.
\end{itemize}

% =====================================================================
% NEW SECTION (skeleton) - Findings: từ quá trình làm thesis,
% từ phản hồi người dùng, và từ 3 chuyên gia tâm lý.
% =====================================================================
\section{Key Findings}
\label{sec:expert-findings}

Beyond the system itself, the design, trial, and consultation process surfaced a number of findings worth stating explicitly. % TODO: câu mở đầu, viết lại cho mượt.

\subsection{Findings from the Design and Development Process}

% TODO: viết lại. Ý chính (rút ra từ Chapter 3 & 4):
\begin{itemize}
\item \textbf{Unifying self-care and therapy is feasible behind one permission model.} A single microservices platform can serve both halves of the care journey without merging the data, as long as sharing is explicit, granular, and revocable.
\item \textbf{Network-level trust between services is not enough.} The most instructive security finding was that an internal caller (the AI service) could reach a user's tracking data with no explicit grant - implicit inter-service trust is exactly the class of assumption the permission model must remove, even for internal callers.
\item \textbf{Grounding the AI in the user's own data, with a safe fallback, is what makes it useful without over-reaching.} % TODO: mở rộng nếu muốn.
\item \textbf{[placeholder]} % TODO: thêm finding kỹ thuật khác nếu cần.
\end{itemize}

\subsection{Findings from User Feedback}

% TODO: viết lại. Cần điền số người dùng thử / phương pháp khảo sát.
\begin{itemize}
\item \textbf{Daily journaling beats periodic surveys.} Trial users found the weekly/monthly surveys redundant next to the daily journal and 3-hourly mood check-ins; the surveys were dropped. Finding: high-frequency, low-friction daily logging yields richer and more meaningful data than periodic questionnaires.
\item \textbf{The most-loved features were the Treasure Box and daily tracking} (Section~\ref{...}), % TODO: cross-ref section trên
which is why they are highlighted separately above.
\item \textbf{[placeholder]} % TODO: điền các phản hồi khác của người dùng thử, số liệu đánh giá, trích dẫn. Mô tả cỡ mẫu và cách khảo sát (N người dùng thử, thời gian).
\end{itemize}

\subsection{Findings from Expert Consultations}

We consulted three mental-health professionals during the project: % TODO: có thể chuyển thành bảng.
\begin{itemize}
\item \textbf{Miss Pham Thi Thanh Qui} - counsellor, Psychological Counselling Office, University of Science, VNU-HCM (in person, 29 May 2026).
\item \textbf{Mister Vuong Nguyen Toan Thien} - CEO, LUMOS (online, 2 June 2026).
\item \textbf{Doctor Pham Tran Thanh Nghiep} - physician, Hoan My Sai Gon Hospital (online, 3 June 2026).
\end{itemize}

\noindent\textbf{What the experts validated and praised.} % TODO: viết lại thành văn.
\begin{itemize}
\item \textbf{Miss Qui} endorsed the option-based emotional journal, the gentle 3-hourly mood check-ins (for complete, meaningful data and for helping users pay attention to themselves), the 4-7-8 breathing exercise as a response to anger, the reward/trophy loop for completing a chain of self-care tasks, the revocable trusted-person sharing, and physical-health tracking such as steps.
\item \textbf{Mister Thien} contributed and praised the Treasure Box concept, the emotion timeline for complete mental-health data, the ``psychology first-aid'' bundle for negative moods (breathing, SOS contact, Treasure Box), and the SOS hotline page.
\item \textbf{Doctor Nghiep} gave the strongest endorsement of the daily reflective tracking features (``already very good'') and suggested adding curated meditation videos to the home screen for users to explore.
\end{itemize}

\noindent\textbf{Cross-cutting findings and design implications.} % TODO: viết lại thành văn.
\begin{itemize}
\item \textbf{Caution about AI in mental health.} All three experts were cautious about AI: two were broadly against relying on it, and one felt the AI should answer only very simple questions, refer harder ones to professionals, and not be given the user's private daily-tracking data. This validates uMatter's conservative AI design - grounded only under an explicit grant, with a safe ungrounded fallback and immediate crisis referral - and motivates the clinical review of the AI companion listed in Future Work.
\item \textbf{Caution about peer support.} Doctor Nghiep cautioned that friends may give advice contrary to professional guidance, suggesting clear in-app disclaimers and steering users toward trusted adults (e.g., parents) rather than peers. Design implication for the social layer. % TODO: nêu rõ đã/ sẽ xử lý thế nào.
\item \textbf{Iterating with experts changed the product.} Concrete features (Treasure Box, dropping weekly surveys, meditation videos, the psychology first-aid flow) came directly from these consultations, not from guesswork - itself a finding about how the app should be built.
\end{itemize}

\section{Limitations}

Three limitations should be kept in mind when interpreting these results. First, this is a software engineering thesis, not a clinical study: no claims are made about the platform's clinical efficacy, and features that would require clinical judgment were explicitly excluded rather than approximated. Second, the security hardening performed was an internal, structured audit rather than an independent third-party penetration test. Third, the deployment evaluated in Chapter~\ref{Chapter4} runs on a single virtual machine and has not been load-tested at a scale representative of real-world production traffic.

\section{Future Work}

Building directly on the limitations above, future work on uMatter falls into three groups:

\begin{enumerate}
\item \textbf{Clinical collaboration.} Any safety-oriented feature that acts on self-reported mental-health data - most notably automated risk detection for at-risk users - requires input from licensed mental-health professionals before it can be responsibly designed, and is intentionally left out of this thesis pending that collaboration. A layered risk-detection framework, informed by clinical guidance, is a natural next step once that collaboration is in place.
\item \textbf{Safety-oriented features more broadly.} Beyond risk detection, this includes guardian-visibility controls appropriate for a teenage user base, clearer in-app pathways to crisis resources, and clinical review of the AI companion's behavior in sensitive conversations.
\item \textbf{Operational maturity.} Independent security review beyond the internal audits performed here, load and scalability testing of the current single-VM deployment, and a multi-instance or managed-database deployment topology to remove the single VM as a point of failure.
\end{enumerate}

Pursued together, these directions would move uMatter from a functioning thesis prototype toward a platform that could be responsibly evaluated with real teenage users under appropriate clinical and institutional oversight.
